\section{Knight's Tour}
\label{sec:cses-knights-tour}

\problemstatement{On an $8\times8$ chessboard, starting from a given square,
find a knight's tour: a sequence of knight moves visiting every square
exactly once. Output the board with each square labelled by the step at
which the knight visits it.}

\begin{patternbox}
A knight's tour is a Hamiltonian path on the knight-move graph. Rather than
exponential search, \textbf{Warnsdorff's heuristic} makes it near-instant:
always move to the square with the \emph{fewest} onward unvisited moves.
\end{patternbox}

\subsection*{Warnsdorff's rule}

From the current square, consider all legal knight moves to unvisited
squares. Move to the one whose \emph{onward degree} -- the number of
unvisited squares \emph{it} can reach -- is smallest. The intuition:
squares with few exits are the easiest to get stranded in, so visit them
early while they are still reachable. Ties can be broken arbitrarily (or by
a fixed order). On a standard board this greedy choice completes a full tour
without backtracking in practice.

\begin{ideabox}[Escape the Tightest Corners First]
Choosing the least-accessible next square keeps the most flexible squares
for later, avoiding dead ends. This ``most-constrained first'' heuristic
turns an exponential Hamiltonian search into an $O(64)$ greedy walk.
\end{ideabox}

\subsection*{Algorithm}

\begin{enumerate}
  \item Mark the start square as step $1$.
  \item Repeat $63$ times: among unvisited knight-move targets, pick the
        one with the fewest unvisited onward moves; go there, label the
        next step.
  \item Output the $8\times8$ grid of step labels.
\end{enumerate}

\subsection*{Dry run}

\dryrun{Start at a corner, e.g.\ $(0,0)$.}

\begin{itemize}
  \item The corner has only two moves; Warnsdorff steers toward low-degree
        squares first, threading all $64$ cells.
  \item The result is a valid tour numbered $1\ldots64$.
\end{itemize}

\subsection*{C++ implementation}

\begin{lstlisting}
#include <bits/stdc++.h>
using namespace std;

int dr[] = {-2,-2,-1,-1,1,1,2,2};
int dc[] = {-1,1,-2,2,-2,2,-1,1};

int main() {
    int sr, sc;
    cin >> sc >> sr;                              // column, row (1-indexed)
    sr--; sc--;

    vector<vector<int>> board(8, vector<int>(8, 0));
    auto inb = [](int r, int c) { return r >= 0 && r < 8 && c >= 0 && c < 8; };
    auto onward = [&](int r, int c) {             // unvisited onward moves
        int cnt = 0;
        for (int k = 0; k < 8; k++) {
            int nr = r + dr[k], nc = c + dc[k];
            if (inb(nr, nc) && board[nr][nc] == 0) cnt++;
        }
        return cnt;
    };

    int r = sr, c = sc;
    board[r][c] = 1;
    for (int step = 2; step <= 64; step++) {
        int best = -1, bestR = -1, bestC = -1;
        for (int k = 0; k < 8; k++) {
            int nr = r + dr[k], nc = c + dc[k];
            if (!inb(nr, nc) || board[nr][nc]) continue;
            int deg = onward(nr, nc);
            if (best == -1 || deg < best) { best = deg; bestR = nr; bestC = nc; }
        }
        r = bestR; c = bestC;
        board[r][c] = step;
    }

    for (int i = 0; i < 8; i++) {
        for (int j = 0; j < 8; j++) cout << board[i][j] << " ";
        cout << "\n";
    }
    return 0;
}
\end{lstlisting}

\begin{complexitybox}
\textbf{Time Complexity.} $O(64)$ moves, each evaluating a constant number
of candidates -- effectively constant for the fixed board. \textbf{Space
Complexity.} $O(1)$ (a fixed $8\times8$ board).
\end{complexitybox}

\begin{takeawaybox}
Warnsdorff's ``fewest onward moves'' heuristic finds a knight's tour without
backtracking by always escaping the most constrained square first. It is the
practical alternative when a full Hamiltonian search would be too slow.
\end{takeawaybox}
